\bookchapter{Глава 56. Шлифовка и публикация — тесты, отладка и выход в Google Play}

\chapterintro{Приложение готово. Что дальше? В этой главе — три дела на пути «от работает» к «можно передавать»: как быстро отлаживать и находить баги, как писать тесты против поломок и как выложить приложение в Google Play, чтобы его качал весь мир. Публикация — не просто красивый ритуал: это водораздел между «самоделкой» и «продуктом».}

\booksection{Три кита отладки: Logcat, точки останова, инспектор UI}

\begin{center}
\begin{tabularx}{\linewidth}{|X|X|X|}
\hline
\textbf{Инструмент} & \textbf{Для чего} & \textbf{Как пользоваться} \tabularnewline
\hline
Logcat & Смотреть системные и свои журналы & В коде Log.d("TAG","значение=\$x"), в панели ищите по TAG, чтобы отфильтровать \tabularnewline
Точки останова & Останавливают программу на нужной строке, чтобы посмотреть значения переменных & Кликните рядом с номером строки — появится красная точка; запустите в режиме Debug и идите по шагам \tabularnewline
Layout Inspector & Смотреть структуру работающего интерфейса & Меню Tools Layout Inspector: дерево компонентов в реальном времени \tabularnewline
\hline
\end{tabularx}
\end{center}

\begin{center}\small\color{gray}Три кита отладки\end{center}

\begin{tcbinfo}{Совет: как читать сообщение об ошибке?}
После падения в Logcat останется красный стек (Stacktrace). Идите снизу вверх и найдите первую строку, начинающуюся с имени вашего пакета (например, com.example.yourapp), — это и есть место преступления. Сначала гляньте «тип Exception» (нулевой указатель? ошибка приведения типа?), потом строку «at ... ваш\_файл:строка». Этот приём закрывает 80\% падений.

\end{tcbinfo}

\booksection{Тесты: страховка для кода}

Ручной проверке верить нельзя: починил одну функцию — тихо сломалась соседняя. Тесты — сразу и «таблетка от сожалений», и «бронежилет». Тестирование в Android трёх уровней; новичку сперва нужно освоить первый:

\begin{center}
\begin{tabularx}{\linewidth}{|X|X|X|X|}
\hline
\textbf{Уровень} & \textbf{Что проверяет} & \textbf{Где лежит} & \textbf{Скорость} \tabularnewline
\hline
Модульные тесты & Чистая логика (вычисления, проверки) & app/src/test/ & Секунды \tabularnewline
Инструментальные тесты & Поведение интерфейса на настоящем устройстве & app/src/androidTest/ & Минуты \tabularnewline
Ручное тестирование & Впечатления живого человека & —— & Нужно всегда \tabularnewline
\hline
\end{tabularx}
\end{center}

\begin{center}\small\color{gray}Три уровня тестов\end{center}

\begin{lstlisting}[language=Kotlin]
// Выносим логику вроде «строка суммы -> число» в чистую функцию
fun parseAmount(input: String): Double? {
    val v = input.trim().toDoubleOrNull()
    return if (v != null && v > 0) v else null
}

// Тестируем (JUnit4)
@Test
fun emptyInputReturnsNull() {
    assertEquals(null, parseAmount(""))
}
@Test
fun negativeReturnsNull() {
    assertEquals(null, parseAmount("-5"))
}
@Test
fun normalValueParses() {
    assertEquals(12.5, parseAmount("12.5"))
}
\end{lstlisting}

Проверку сумм и статистические расчёты в учётной книге обязаны прикрывать модульные тесты. Запуск: правый клик по папке test Run. Зелёный — всё прошло, меняйте код спокойно.

\booksection{Подпись и сборка: ставим на приложение «печать»}

Android требует подписи для всех устанавливаемых пакетов — как печать на документе: она подтверждает, что «пакет выпущен мной и не подделан». При отладке Android Studio автоматически подписывает debug-подписью; для выпуска нужна собственная release-подпись. Порядок:

\begin{enumerate}
  \item меню Build Generate Signed Bundle / APK;
  \item создайте файл keystore (хранилище ключей) — задайте пароль и храните бережно;
  \item выберите release и формат AAB — сборка готова.
\end{enumerate}

\textbf{Осторожно: потеряли keystore — потеряли приложение}\par

После публикации обновления подписываются тем же ключом, иначе магазин отклонит обновление. Потеряли файл keystore и пароль — останется лишь выпустить приложение заново как новое, и старые пользователи останутся ни с чем. Обязательно сделайте резервную копию keystore и пароля в надёжном месте (зашифрованное облако + блокнот) — это один из самых дорогих уроков Android-разработчика.

\booksection{Публикация в Google Play: весь путь}

Зарегистрировать аккаунт разработчика

Подготовить материалы для магазина

Загрузить AAB

Заполнить форму безопасности данных

Проверка (от часов до дней)

Публикация

\begin{center}\small\color{gray}Рисунок 6. Шесть шагов публикации в Google Play\end{center}

\begin{itemize}
  \item Регистрация: разовый регистрационный взнос 25 долларов; нужен работающий Google-аккаунт;
  \item Материалы: значок приложения (512×512), скриншоты (минимум 2), описание функций, ссылка на политику конфиденциальности (обязательна, если собираете данные);
  \item Форма безопасности данных: честно перечислите, что собираете, — учётная книга без сети может указать «сбор отсутствует»;
  \item Проверка: первую публикацию нового аккаунта часто смотрит живой модератор — от нескольких часов до нескольких дней; вернули — правьте по замечаниям.
\end{itemize}

Пользователям из Китая: для Google Play нужен особый сетевой доступ, а в материковом Китае распространение обычно идёт через местные магазины (Huawei, Xiaomi, Yingyongbao и др.). Процедура в местных магазинах похожа, но у каждого свои вкусы проверки — один и тот же AAB/APK можно разослать в несколько магазинов. Для тренировки достаточно пройти хоть один.

\booksection{Материалы для магазина: первое впечатление — это и есть загрузки}

Страница в магазине — витрина, которую смотрит «зашедший» пользователь. Первое впечатление делают три материала:

\begin{itemize}
  \item Значок: 512×512; не отделывайтесь скриншотом. Нет дизайнерских навыков — сгенерируйте лаконичную геометрическую фигуру ИИ-инструментом; зачёт, только если значок узнаётся с одного взгляда даже уменьшенным на рабочем столе телефона;
  \item Скриншоты: минимум 2, лучше 4–6 — показывайте самые главные сценарии; к каждому добавьте одну простую человеческую фразу («записать трату за секунду» в десять раз лучше «эффективного решения для учёта финансов»);
  \item Описание: первая фраза должна ясно отвечать «что это и для кого»; не лепите горой «самое профессиональное, самое мощное» — поисковый вес в магазине зависит в основном от ключевых слов и конверсии в загрузки, искренность действеннее преувеличений.
\end{itemize}

\begin{tcbinfo}{Совет: сравните себя с соседями по полке}
До публикации прогуляйтесь по первой десятке похожих приложений в магазине и сравните скриншоты: какой у них стиль значка? какой довод вынесен на скриншоты? как написаны заголовки? Подражать — не значит красть: переймите ход, вложите своё содержание. Это самый быстрый урок холодного старта для любого продукта.

\end{tcbinfo}

\booksection{Мини-чек-лист: производительность и отчёты о сбоях}

Перед публикацией приложения на эти два дела стоит потратить по полчаса:

\booksubsection{Четыре пункта проверки производительности}

\begin{enumerate}
  \item Списки — через LazyColumn: длинные списки ни в коем случае не собирайте жёстким Column — ленивая загрузка это первый урок производительности Android;
  \item Не пихайте картинки в оригинальном размере: библиотека Coil сама уменьшит и закэширует; один оригинал на 5 МБ способен задушить слабый телефон;
  \item Не гоняйте тяжёлую работу в главном потоке: база данных, сеть, большие файлы — всё в корутинах;
  \item Измеряйте через Profiler: встроенный в Android Studio анализатор производительности показывает память, CPU и расход батареи — сначала измерьте, потом оптимизируйте, на ощущения не надейтесь.
\end{enumerate}

\booksubsection{Отчёты о сбоях}

То, что у вас при тестах не падало, не значит, что не упадёт у пользователя. До публикации подключите Firebase Crashlytics (официальный сервис Google, бесплатный): несколько строк настроек — и при падении стек сам прилетит в консоль: увидите «какая модель, какая версия, на какой строке рухнуло». Публиковать без мониторинга сбоев — всё равно что ездить с закрытыми глазами.

\begin{tcbpractice}{Практика: «сдайте» учётную книгу в эксплуатацию}
(1) Добавьте 3 модульных теста к логике вроде parseAmount; (2) через Logcat своими руками поймайте специально подстроенное падение и почините его; (3) соберите APK с release-подписью и отправьте другу на установку (вот она, церемония «произведения»); (4) если есть силы — зарегистрируйте аккаунт разработчика и действительно отправьте учётную книгу на публикацию. Первые три пункта по силам уже сегодня вечером.

\end{tcbpractice}

\begin{tcbsummary}{Коротко о главном}
\begin{itemize}
  \item Три кита отладки: Logcat — журналы, точки останова — переменные, Layout Inspector — интерфейс
  \item Читать падение: в стеке снизу вверх ищите первую строку с именем своего пакета
  \item Три уровня тестов: модульные (секунды, осваиваем первыми) / инструментальные / ручные
  \item Выпуск — только с подписью: держите резервные копии keystore и пароля; потеряете — приложение «сменит хозяина»
  \item Шесть шагов публикации: регистрация материалы загрузка AAB форма безопасности данных проверка публикация
\end{itemize}

\end{tcbsummary}
